\documentclass[11pt]{article}
\usepackage[T1]{fontenc}
\usepackage[utf8]{inputenc}
\usepackage{lmodern,amsmath,amssymb,amsthm,mathtools}
\usepackage[margin=1in]{geometry}
\usepackage{microtype}
\usepackage[colorlinks=true,linkcolor=blue,citecolor=blue,urlcolor=blue]{hyperref}
\newtheorem{theorem}{Theorem}[section]
\newtheorem{proposition}[theorem]{Proposition}
\newtheorem{lemma}[theorem]{Lemma}
\newtheorem{corollary}[theorem]{Corollary}
\theoremstyle{definition}\newtheorem{definition}[theorem]{Definition}
\theoremstyle{remark}\newtheorem{remark}[theorem]{Remark}
\newcommand{\dd}{\mathrm d}
\newcommand{\vol}{\dd V}
\newcommand{\supp}{\operatorname{supp}}
\newcommand{\Hess}{\operatorname{Hess}}
\newcommand{\Div}{\operatorname{div}}
\newcommand{\tr}{\operatorname{tr}}
\newcommand{\Op}{\operatorname{Op}}
\newcommand{\ip}[2]{\langle #1,#2\rangle}
\newcommand{\norm}[1]{\lVert #1\rVert}
\title{Uniqueness of a smooth real potential on a known Riemannian manifold}
\author{Kaiwen Jin\\\small\href{mailto:yui.ui.siki@gmail.com}{yui.ui.siki@gmail.com}}
\date{October 13, 2026}
\begin{document}
\maketitle
\begin{abstract}
We prove uniqueness of a smooth real scalar potential from its full
Dirichlet-to-Neumann map at one energy on an arbitrary known compact,
connected, smooth Riemannian manifold of dimension at least three with
smooth nonempty boundary. The only spectral assumption is that this energy
is not a Dirichlet eigenvalue; negative eigenvalues and sign-changing
potentials are allowed. Equal boundary maps first give common exterior
source solutions through a weak gluing argument that does not require a
positive quadratic form. Near an interior frontier, a common source gives
local positive conjugations of the two equations. We adapt the separately
elliptic continuation and shrinking-sphere method of OpenAI to the resulting
operators with the same principal metric. An angular coercivity estimate
bounds the actual transfer densities jointly in center and direction. Their
collapse gives equality on one neighborhood for every matching source pair.
Finite source jets then recover the local conjugating solutions and the
potentials, closing the frontier without a global positive solution.
\end{abstract}

\section{Introduction}\label{sec:intro}

Let $(M,g)$ be a known compact connected smooth Riemannian manifold of
dimension $n\ge3$ with smooth nonempty boundary. For a real smooth
potential $q$ such that zero is not a Dirichlet eigenvalue of
$L_q=-\Delta_g+q$, define the ordinary Dirichlet-to-Neumann map by
\begin{equation}\label{eq:DN}
 \Lambda_{g,q}f=\partial_{\nu_g}u|_{\partial M},\qquad
 L_qu=0\ \text{in }M,\qquad u|_{\partial M}=f,
\end{equation}
where $\nu_g$ is the outward unit normal. The map acts from
$H^{1/2}(\partial M)$ to $H^{-1/2}(\partial M)$; equivalently, its
values and pairings on all smooth boundary data determine the map.
The boundary is labeled, and both potentials are compared on the same
manifold with the same metric. We ask whether this single boundary
response determines the interior scalar coefficient without a geometric
restriction on $g$.

The fixed-energy question must be distinguished from recovery using a
family of energies or a sufficiently large frequency together with a ray
transform hypothesis. For example, Ma, Sahoo and Salo
\cite{MaSahooSalo} establish a large-fixed-frequency result on manifolds
with an invertible ray transform. Such results do not give the assertion
below on an arbitrary known smooth manifold at the prescribed energy zero.
We do not assume simplicity, absence of conjugate points, a limiting
Carleman weight, or injectivity of a geodesic ray transform.

Our continuation construction adapts the mixed-kernel transfer, product
continuation, shrinking-sphere construction, and uniform density and
coercivity estimates of OpenAI \cite{OpenAIMetric}. A separately elliptic
mixed Green kernel is continued along product hypersurfaces, and its
transfer integrals on shrinking spheres are controlled through angular
raising and lowering operators. The later metric and unitary-connection preprint
\cite{OpenAIConnection} addresses an operator of the form $d_A^*d_A$;
its theorem does not include an additional independent real scalar
potential. We give the continuation and angular arguments needed here,
with their overlap and parameter dependence specified, and adapt them to
two locally conjugated scalar equations having the same principal metric.
The core continuation method is credited to these sources.

\begin{theorem}\label{thm:main}
Let $(M,g)$ be any compact connected smooth Riemannian manifold of
dimension $n\ge3$ with smooth nonempty boundary. Let
$q_1,q_2\in C^\infty(M;\mathbb R)$, and suppose zero is not a Dirichlet
eigenvalue of either $L_j=-\Delta_g+q_j$, with Dirichlet domain
$H^2(M)\cap H^1_0(M)$. If
\[
 \Lambda_{g,q_1}=\Lambda_{g,q_2}
\]
as full ordinary boundary maps at energy zero, then $q_1=q_2$ on $M$.
\end{theorem}

No orientability or connectedness of the boundary is required. The metric
is known throughout, and no diffeomorphism or coefficient gauge occurs in
the conclusion. Both potentials are arbitrary real smooth functions
subject to the stated nonresonance condition. In particular, either
Dirichlet operator may have negative eigenvalues, and the difference of
the potentials may be flat or change sign arbitrarily often.

The spectral qualification matters for the proof. A global positive
solution $s$ of $(-\Delta_g+q)s=0$ would give, for interior compactly
supported $u$,
\[
 \int_M\bigl(|\dd u|_g^2+q|u|^2\bigr)\vol_g
       =\int_M s^2|\dd(u/s)|_g^2\vol_g\ge0.
\]
By density this excludes negative Dirichlet eigenvalues. Thus a global
positive conjugation would impose a stronger assumption than the theorem.
Shifting the operator to obtain positivity would also change the measured
energy. We instead use nonresonance for the signed weak inverse and select
positive source solutions only in a small neighborhood of a frontier.

Here is the structure of the argument. Equality of the boundary maps gives
equal boundary jets of the potentials. A thin exterior cap can therefore
be attached with common coefficients. Nonresonance persists for a small
cap, and weak stationarity transfers every source supported in the cap
between the two Dirichlet inverses. Point distributions and scalar unique
continuation show that these sources have arbitrary local value and first
jets subject only to their differential equation.

At a point where all source families agree from one side, one matching
source is positive nearby. Dividing each equation by its own positive
solution removes the zeroth-order term locally. A compatible mixed kernel
then defines a transfer density on small spheres. Strict product
hypersurfaces continue this kernel at every positive radius; no uniform
kernel bound at radius zero is asserted. The complete radial map, including
its smooth remainder, gives an equation for the actual density. Concavity
of a reference depth function supplies an angular coercivity estimate and
a uniform joint $L^2$ bound in the center and sphere variables.

The radii collapse on one fixed neighborhood. The joint bound makes every
matching solution pair agree there, with that neighborhood chosen before
the further source. A finite spanning set of source gradients then
identifies the two local positive conjugations and the original potentials.
This common-neighborhood conclusion advances the open set of agreement
past any putative interior frontier and proves the theorem globally.

Section~\ref{sec:cap} establishes the signed inverse, cap data and source
jets. Sections~\ref{sec:local}--\ref{sec:spheres} construct the local
conjugation and the continued mixed kernel. Sections~\ref{sec:trace}
and~\ref{sec:coercivity} derive and bound the transfer density.
Section~\ref{sec:global} closes the global frontier. We use
$\Delta_g=\Div_g\nabla_g$, positive volume densities, and the outward
normal convention in \eqref{eq:DN}. All coefficients are smooth up to
any boundary on which they are used. The proof does not assert a partial
boundary-data theorem, a complex-potential theorem, or a resonant case.


\section{Signed Dirichlet inverses and common exterior data}\label{sec:cap}

We first obtain interior source data from the boundary response. This step
must allow an indefinite Dirichlet form. Throughout, $\Delta_g=\Div_g\nabla_g$,
$L_j=-\Delta_g+q_j$, and all integrals on manifolds use positive densities.
For real functions put
\[
 \mathfrak b_j(u,v)=\int_M(\ip{\dd u}{\dd v}_g+q_juv)\vol_g.
\]
The complex extension is sesquilinear, linear in its first variable. The
Dirichlet realization has domain $H^2(M)\cap H^1_0(M)$, where $H^1_0$ is
the closure of interior compactly supported smooth functions. Smooth data
suffice for every boundary pairing below.

\begin{lemma}\label{lem:inverse}
If $0$ is not a Dirichlet eigenvalue of $L=-\Delta_g+q$ with real smooth
$q$, its weak realization $H^1_0\to H^{-1}$ is a bounded isomorphism.
This property persists under a sufficiently small smooth outward
perturbation of the boundary and coefficients, pulled back to $M$.
\end{lemma}
\begin{proof}
Add a constant $c>0$ so that $\mathfrak b_c$ is coercive on $H^1_0$.
Lax--Milgram gives its inverse $R_c$ \cite[Theorem 4.20]{HunterPDE}.
The inclusion $H^1_0\hookrightarrow L^2$ is compact: extend its defining
approximants by zero in boundary charts, localize in finitely many cubes,
and truncate Fourier frequencies. The $H^1$ bound makes the remaining
$L^2$ norm $O(N^{-1})$. Thus $K=cR_cJ$, with $J$ the $L^2$ pairing, is
compact and self-adjoint in the $\mathfrak b_c$ inner product.
A weak vector in $\ker(I-K)$ gains $H^2$ by boundary elliptic regularity
and is a zero Dirichlet eigenfunction; hence this kernel is zero.
If $I-K$ were not bounded below, a unit sequence and compactness of $K$
would produce a nonzero kernel vector. Its range is therefore closed;
its orthogonal complement is $\ker(I-K)^*=0$, so it is onto.
This proves the asserted inverse. Boundary regularity here is the
pure divergence estimate of \cite[Theorems 4.30--4.31]{HunterPDE}; smooth
lower terms enter the forcing and are bootstrapped in finite charts.

After pullback, an outward perturbation changes the weak form in operator
norm by a quantity tending to zero. If $R$ is the original inverse,
$\norm{R(\mathfrak b_\epsilon-\mathfrak b_0)}<1$ gives the perturbed inverse
by a Neumann series. Positivity of the original form is unnecessary.
\end{proof}

Equality of the DN maps determines the complete boundary jets of $q_1,q_2$.
For the identical metric this is the scalar, zero-connection case of the
last clause of \cite[Theorem 3.4]{Cekic}; the normal gauges are identity.
To fix the sign, in inward boundary distance $s$ let
$\rho=(g^{ab}\xi_a\xi_b)^{1/2}$. The factor symbol starts with $b_1=-\rho$,
and $b_0$ depends only on $g$. Its scalar recursion gives
\begin{equation}\label{eq:jets}
 b_{-k-1,1}-b_{-k-1,2}
 =-\frac{\partial_s^k(q_1-q_2)}{(2\rho)^{k+1}}
   +\text{terms fixed by lower potential jets and metric jets},\qquad k\ge0.
\end{equation}
Equal complete DN symbols recover these jets inductively. The simultaneous
metric normalization in the first clause of that theorem is not needed
when the metric is already identical. Flatness at the boundary gives jets,
not equality on an interior collar.

Choose a boundary patch and a nonnegative smooth bump with connected
positivity set. Attach a thin cap $E$ across that patch, obtaining a smooth
compact manifold $N$. Extend the metric smoothly and use one common exterior
extension for the two potentials. This is possible by \eqref{eq:jets}:
their difference has a smooth zero extension. A smooth extension of either
coefficient can be constructed from cutoff Taylor monomials, choosing the
$j$th support small enough that its derivatives of order less than $j$
are bounded by $2^{-j}$. The series converges in every fixed smooth norm.
A collar diffeomorphism identifies $N$ with a small boundary perturbation
of $M$. Lemma~\ref{lem:inverse} makes both extended Dirichlet problems
invertible after the cap is made sufficiently thin.

Write $R_j=L_{j,D}^{-1}$ on $N$, and normalize its kernel by
\[
 R_jF(x)=\int_N G_j(x,y)F(y)\vol_g(y).
\]
The kernels are real and symmetric. An interior order $-2$ parametrix $Q$
with $L_jQ=I+S$ gives $G_j=Q-R_jS$ locally; $S$ and $R_jS$ have smooth
kernels. The off-diagonal symbol-kernel estimate
\cite[Lecture I, (1.21)--(1.24)]{Taylor} gives
\begin{equation}\label{eq:kernelbound}
 |\partial_x^a\partial_y^bG_j(x,y)|
 \le C_{a,b}|x-y|^{2-n-|a|-|b|}
\end{equation}
on sufficiently small compact coordinate neighborhoods, with a smooth
remainder absorbed into the constant. Its order condition is satisfied
because $n\ge3$.

\begin{lemma}\label{lem:cap}
For every $F\in C_c^\infty(E)$, $R_1F=R_2F$ on $E$.
Consequently $G_1=G_2$ on $E\times E$.
\end{lemma}
\begin{proof}
Let $w_1=R_1F$, and replace its restriction to $M$ by the $L_2$ solution
with the same boundary trace $h$. Leave it equal to $w_1$ on $E$.
The replacement differs on $M$ by a member of $H^1_0(M)$; its zero
extension gives a glued member of $H^1_0(N)$. In the weak equation,
the difference of the interior boundary pairings is
$\ip{(\Lambda_2-\Lambda_1)h}{\phi|_{\partial M}}=0$.
The cap coefficients and forcing are common, so the glued function solves
$L_2w=F$. Uniqueness identifies it with $R_2F$.
This is a stationarity argument for the weak form, not minimization.
Testing all cap sources and using smoothness off the pole gives the kernel
identity; the equal delta singularities also cancel distributionally.
\end{proof}

Fix a nonempty source set $U_0\Subset E$.
\begin{lemma}\label{lem:jets}
At any $p\in N^\circ\setminus\overline{U_0}$, the second jets of
$R_jF$, $F\in C_c^\infty(U_0;\mathbb R)$, fill exactly the hyperplane
$L_jv(p)=0$. In particular, value and first derivatives are arbitrary.
\end{lemma}
\begin{proof}
Let an order-at-most-two point distribution $T$ annihilate these jets.
Define $R_jT$ by transposition. Then $R_jT=0$ on $U_0$ and
$L_jR_jT=T$. Away from $p$ it is smooth by a local parametrix and vanishes
by scalar unique continuation \cite{Aronszajn}. The punctured interior
is connected in dimension at least two: paths can detour around a point
in a coordinate ball. Thus $R_jT$ is supported at $p$.

A point-supported distribution is a finite sum of delta derivatives.
Indeed, a test whose derivatives through its finite order vanish at $p$
can be cut off at radius $\epsilon$; Taylor's remainder makes its test
seminorm $O(\epsilon)$ and its pairing vanish. For a nonzero such
distribution of order $k$, an elliptic second-order operator has exact
image order $k+2$: its top polynomial is multiplied by a nonzero
quadratic principal symbol. Hence $R_jT=c\delta_p$ and $T=cL_j\delta_p$,
where distribution transposes use $\vol_g$. This is precisely the
annihilator of the equation hyperplane. Finite-dimensional duality proves
surjectivity onto it. The Hessian trace can enforce the equation while
value and first derivatives are prescribed. Real parts give real sources.
\end{proof}

\section{Local positive conjugation}\label{sec:local}

At a point where the two source families agree from one side, choose a
matching pair with common value one using Lemma~\ref{lem:jets}. On a fixed
small neighborhood it gives positive solutions $s_j$ of $L_js_j=0$.
For $u_j=s_jv_j$ the exact conjugation is
\begin{equation}\label{eq:ground}
 s_j^{-1}L_j(s_jv_j)
 =-\Delta_gv_j-2\ip{\dd\log s_j}{\dd v_j}_g
 =:A_jv_j=-s_j^{-2}\Div_g(s_j^2\nabla_gv_j).
\end{equation}
Thus $A_j1=0$ and both principal metrics are identical. Their local
symmetric densities are $s_j^2\vol_g$, and
\[
 K_j(x,y)=\frac{G_j(x,y)}{s_j(x)s_j(y)}
\]
is the corresponding local Green kernel. Only local positivity is needed.

The following statement is the local result used in the global proof.
The kernels in it are the actual kernels just constructed on $N$.
\begin{proposition}\label{prop:local}
Suppose a neighborhood $\Omega$ of $0$ has a known open subset $W$ containing
one side of a smooth hypersurface through $0$. Suppose $s_1,s_2>0$ solve
$L_js_j=0$ on $\Omega$, agree on $W$, and $G_1=G_2$ on $W\times W$.
There is one neighborhood $Z$ of $0$ such that every pair of smooth
solutions $A_jv_j=0$ on $\Omega$ agreeing on $W$ satisfies $v_1=v_2$ on $Z$.
\end{proposition}

We prove this through a transfer integral on shrinking spheres. The
neighborhood $Z$ and the integral density will depend on the fixed kernels,
groundings and geometry, but not on a further solution pair.
Choose coordinates with $g(0)=I$ and known side tangent to $x_n<0$.
Under $x\mapsto\lambda x$, set
\[
 g_\lambda(x)=g(\lambda x),\quad s_{j,\lambda}(x)=s_j(\lambda x),\quad
 K_{j,\lambda}(x,y)=\lambda^{n-2}K_j(\lambda x,\lambda y).
\]
On fixed bounded sets the principal metric tends smoothly to Euclidean,
and the grounded drift tends to zero. Each compact subset strictly below
the plane becomes known for small $\lambda$. Estimate~\eqref{eq:kernelbound}
gives uniform smooth bounds on compact separated pairs. Initially use
$K_1$ when $y$ is known and $K_2$ when $x$ is known. They agree on the
intersection and satisfy $A_{1,x}K=A_{2,y}K=0$ off the diagonal.
For the rest of the local proof subscripts $\lambda$ are suppressed.
The small spatial scale will be fixed before any radius tends to zero.

\section{Separately elliptic continuation}\label{sec:cross}

We give the continuation argument since its domain and overlap properties
are used below. Scalar local regularity, quantitative interior estimates,
and the calculus used here follow from an elliptic parametrix
\cite[Lecture I, Sections 3--4]{Taylor}. On compact buffers the required
integer estimates have finite coefficient dependence. For example, testing
$Pu=\partial_iF^i+h$ against $\eta^2\bar u$ gives
\[
 \norm{\nabla u}_{L^2(U)}\le
 C(\norm{u}_{L^2(V)}+\norm{F}_{L^2(V)}+\norm{h}_{L^2(V)}),\qquad U\Subset V.
\]
Differentiating a homogeneous equation gives divergence forcing involving
only already controlled derivatives, and finite nested buffers give every
integer order. Initially distributional $L^2$ solutions gain local Sobolev
regularity by the parametrix, so the weak tests are legitimate. Constants
use only finitely many smooth coefficient bounds at each specified order.

\subsection{Strict scalar weights}
Let $\theta$ have nonzero differential and, with $Z=\nabla\theta/|\dd\theta|$,
suppose
\begin{equation}\label{eq:strictscalar}
 \Hess\theta(Z,Z)+\Hess\theta(V,V)>0\quad(V\perp Z,\ |V|=1).
\end{equation}
On a small neighborhood choose $m$ strictly between minus the least
tangential Hessian eigenvalue and the normal Hessian value. For
$v=e^{\tau\theta}w$, the conjugated positive Laplacian is
\[
 A_\tau+B_\tau=-\Delta-\tau^2|\dd\theta|^2
       +\tau(2\nabla\theta\cdot\nabla+\Delta\theta),
\]
where $A_\tau$ is self-adjoint and $B_\tau$ skew-adjoint on compact supports.
Integration by parts gives
\begin{align*}
 \norm{(A_\tau+B_\tau)v}^2
 &=\norm{A_\tau v}^2+\norm{B_\tau v}^2+\ip{[A_\tau,B_\tau]v}{v},\\
 \ip{[A_\tau,B_\tau]v}{v}
 &=4\tau\int\Hess\theta(\nabla v,\nabla\bar v)
   +4\tau^3\int\Hess\theta(\nabla\theta,\nabla\theta)|v|^2
   -\tau\int\Delta^2\theta|v|^2.
\end{align*}
Ricci terms from derivative commutation cancel those in
$\Div\Hess\theta=\dd\Delta\theta+\operatorname{Ric}(\nabla\theta,\cdot)$.
Add $4m\tau\operatorname{Re}(A_\tau v,v)$ to this form. It makes the
tangential gradient and leading scalar terms positive. The remaining
normal-gradient loss is controlled by
$\norm{\nabla_Zv}^2\le C\tau^{-2}\norm{B_\tau v}^2+C\norm{v}^2$.
The added form costs at most $\frac14\norm{A_\tau v}^2+C\tau^2\norm{v}^2$.
For large $\tau$ these losses are absorbed. Smooth lower terms have
conjugated norm squared at most $C(\norm{\nabla v}^2+\tau^2\norm{v}^2)$
and are absorbed too. For the actual scalar operator $A$ this yields
\begin{equation}\label{eq:carleman}
 \tau\norm{e^{\tau\theta}\dd w}^2+\tau^3\norm{e^{\tau\theta}w}^2
 \le C\norm{e^{\tau\theta}Aw}^2,\qquad \tau\ge\tau_0.
\end{equation}
The constants are uniform on fixed strict margins and finite coefficient bounds.

Suppose a cutoff equals one on $O\Subset D$, its bottom derivative support
has weight at most $a$, the remaining derivative support at most $b$, and
$\theta\ge c$ on $O$, where $a<c<b$. Enlarge the latter support to a compact
observed open set $S\Subset D$. Interior estimates on the enlarged collars
and \eqref{eq:carleman} give, for $Au=0$,
\[
 \norm{u}_O\le C\{e^{-\tau(c-a)}\norm{u}_D+
                       e^{\tau(b-c)}\norm{u}_S\}.
\]
Balancing the terms when the balancing $\tau$ exceeds $\tau_0$, and using
the trivial bound otherwise, proves
\begin{equation}\label{eq:interp}
 \norm{u}_O\le C\norm{u}_S^\gamma\norm{u}_D^{1-\gamma},
 \qquad\gamma=\frac{c-a}{b-a}.
\end{equation}
If the observed norm is zero, let $\tau\to\infty$.

\subsection{Extension from a cross}
\begin{lemma}\label{lem:cross}
Let $D_i$ be bounded connected coordinate domains and $S_i,O_i\Subset D_i$
open, with $S_i$ nonempty. Let $A_i$ be smooth real scalar elliptic
operators with positive principal matrices. Suppose their $L^2$ solutions
satisfy \eqref{eq:interp} with exponents $0<\gamma_i<1$ and
$\alpha=\gamma_1+\gamma_2-1>0$.
Compatible separately homogeneous $L^2$ functions on
$(D_1\times S_2)\cup(S_1\times D_2)$ extend to a smooth separately
homogeneous function on $O_1\times O_2$, agreeing on its actual
intersections with both given pieces. On compact output subsets its
smooth norms are bounded by the summed input $L^2$ norms.
\end{lemma}
\begin{proof}
Let $\mathcal H=\{w\in L^2(D_1):A_1w=0\}$, a closed Hilbert subspace.
Restriction $R:\mathcal H\to L^2(S_1)$ is compact by cutoff Fourier
truncation and injective by unique continuation. On each fixed cutoff
coordinate cube, identify opposite faces and use the torus Fourier basis
\cite[Section 2, Theorem 21 and Corollary 22]{TaoFourier}.
Parseval and the integer derivative multipliers give, for each integer $m$,
a rank-$O(N^{n_1})$ approximation with norm error $O(N^{-m})$,
where $n_1=\dim D_1$.

Here the compact spectral construction can be made directly. Maximize
$\norm{Rw}^2$ on the unit sphere using a weakly convergent maximizing
sequence and compactness of $R$. The maximizing vector has norm one;
variation in orthogonal directions makes its restriction orthogonal to
all subsequent ones. Iterate on orthogonal complements. The resulting
basis $w_j$ is complete: a vector perpendicular to all of it has image
norm bounded by every decreasing singular value and hence zero.
In the infinite case compactness forces these values to zero.
Thus
\[
 \norm{w_j}_{D_1}=1,\quad (Rw_i,Rw_j)_{S_1}=\mu_j^2\delta_{ij},\quad
 0<\mu_j\le1.
\]
A vector in the span of the first $j$ basis vectors annihilated by a
rank-less-than-$j$ approximation has image norm at least $\mu_j$.
The Fourier error therefore gives $\mu_j\le C_mj^{-m/n_1}$, and in
particular decay faster than any fixed power. No lower bound is used.
The same argument covers finite-dimensional spaces.

Write the first cross piece $F_1=\sum_jw_jb_j^0$ on $D_1\times S_2$.
For the second piece define
\[
 b_j(y)=\mu_j^{-2}\int_{S_1}F_2(x,y)\overline{w_j(x)}\dd x.
\]
If the summed input norm is $M$, Cauchy--Schwarz gives
$\norm{b_j}_{D_2}\le M/\mu_j$ and compatibility gives
$b_j=b_j^0$ on $S_2$, $\norm{b_j}_{S_2}\le M$.
Fubini and compact-test approximation in $x$ give $A_2b_j=0$.
Interpolation then gives
\[
 \norm{w_j}_{O_1}\le C\mu_j^{\gamma_1},\quad
 \norm{b_j}_{O_2}\le CM\mu_j^{\gamma_2-1},\quad
 \norm{w_jb_j}_{O_1\times O_2}\le CM\mu_j^\alpha.
\]
The series converges absolutely in product $L^2$ and solves both equations.
Their sum is elliptic on the product, giving smoothness and compact-subset
estimates.

To identify the second piece, put $V=\overline{R\mathcal H}$.
For $z\perp V$, the $A_2$ solution
$\int_{S_1}F_2(x,y)\bar z(x)\dd x$ vanishes on $S_2$ and hence on $D_2$.
A countable dense set of such $z$ proves that almost every second-piece
slice belongs to $V$. Parseval in the basis $Rw_j/\mu_j$ identifies its
coefficients as $\mu_jb_j$. The same partial sums converge to $F_2$ on
$S_1\times D_2$, and absolutely on $S_1\times O_2$ by the bound
$CM\mu_j^{\gamma_2}$. They agree with the product limit on
$(S_1\cap O_1)\times O_2$. Completeness in $\mathcal H$ similarly gives
agreement with $F_1$ on $O_1\times(S_2\cap O_2)$.
For a complete orthonormal basis, Pythagoras and the limit of orthogonal
projections give Parseval. Its squared tails pass through the finite
product integrals by dominated convergence, bounded by the original
squared $L^2$ norm. The interchanges use dominated convergence and
absolutely integrable Fubini
\cite[Theorems 18 and 23]{HunterMeasure}.
If $\mathcal H=0$, the first piece is zero and the same pairings plus
unique continuation make the second piece zero.

For uniform families choose a fixed positive lower bound for $\alpha$,
fixed buffers, uniform interpolation constants and finite coefficient
bounds. One Fourier decay order then makes $\sum_j\mu_j^\alpha$ uniformly
bounded. Smooth output estimates are uniform under these same bounds.
\end{proof}

\subsection{Crossing a product hypersurface}
Let $\rho=0$ be a smooth product hypersurface with both partial gradients
nonzero. Put $a_i=\nabla_i\rho/|\dd_i\rho|^2$ and $H=\Hess\rho$ in the
product principal metric. A sufficient strict condition is
\begin{equation}\label{eq:productcriterion}
 H((a_1,-a_2),(a_1,-a_2))+H((v_1,v_2),(v_1,v_2))>0,
 \quad v_i\perp a_i,\quad |v_i|=|a_i|.
\end{equation}
\begin{lemma}\label{lem:surface}
Under \eqref{eq:productcriterion}, a smooth separately homogeneous function
on $\rho>0$ extends across the point to a product neighborhood, agreeing
with its original germ on the positive side.
\end{lemma}
\begin{proof}
Let $\mathcal C_i$ consist of symmetric forms $P_i$ with
$P_i(a_i,a_i)+P_i(v_i,v_i)>0$ for every indicated $v_i$.
Its dual consists of $d_ia_ia_i^T+S_i$, with $d_i\ge0$,
$S_i\ge0$ tangential and $\tr S_i=d_i|a_i|^2$.
The dual of $(\mathcal C_1\oplus\mathcal C_2)+\mathrm{PD}$ consists of
positive semidefinite matrices with these diagonal blocks.
Its trace-one slice is compact. If a dual matrix $Q$ annihilates
$\dd\rho$, covariance decomposition makes the two normal components
opposite, with common variance $d>0$ and no covariance with tangential
components. The tangential covariance has two fixed traces.
Every such covariance is a finite convex combination of rank-one ones:
at rank $r\ge2$, a supported symmetric variation preserving the two traces
exists since $r(r+1)/2>2$. Moving in both directions to the first singular
endpoints lowers rank and expresses the original as their convex
combination; repeat. The rank-one covariances are
$d(v_1,v_2)(v_1,v_2)^T$. Thus \eqref{eq:productcriterion} is positivity of
$\tr(QH)$ on all annihilating dual tests.

Compactness gives a positive margin on those tests. Near them it persists;
away from them $\tr(Q\dd\rho\otimes\dd\rho)$ has a positive lower bound.
Choose $K$ large. Finite-dimensional convex separation then gives
\[
 H+K\dd\rho\otimes\dd\rho>P_1\oplus P_2,\qquad P_i\in\mathcal C_i.
\]
The separation used here follows from nearest-point projection onto a
closed convex set, followed by a limit of supporting normals at a boundary
point. Choose factor functions $\theta_i$ with first jets $\dd_i\rho$ and
Hessians $P_i$, and value zero at contact. Taylor's formula gives
\begin{equation}\label{eq:quadraticgap}
 \theta_1+\theta_2\le\rho+K\rho^2/2-c(|x|^2+|y|^2)
\end{equation}
on a small product neighborhood.

Take factor cylinders of lateral radius $R$ and height $(-l,l+4e)$,
with $l\ll cR^2$ and $e\ll\kappa l^2$. A cutoff is one on the inner
lateral cylinder for $-l+3e<\theta_i<l+2e$; its top derivatives lie in
$l+2e<\theta_i<l+3e$. Observe the enlarged band
$l+e<\theta_i<l+7e/2$ and an enlarged lateral collar.
The top observed cross has height sum greater than $e$; the lateral
cross has quadratic gap larger than the possible negative height sum.
By \eqref{eq:quadraticgap} both crosses lie strictly on $\rho>0$.

Use the strict scalar weights $\phi(\theta_i)=\theta_i-\kappa\theta_i^2$.
Choose $\kappa>0$ small to preserve \eqref{eq:strictscalar}. The bottom,
top derivative and target weights can be taken as
\[
 a=\phi(-l+3e),\quad b=\phi(l+3e),\quad c=\phi(-e).
\]
The output has inner lateral radius and $-e<\theta_i<l+2e$.
Equation~\eqref{eq:interp} gives
$\gamma_i=(c-a)/(b-a)>1/2$ for small $e>0$, since at $e=0$ it equals
$1/2+\kappa l/2$. Lemma~\ref{lem:cross} therefore applies.
The enlarged observed band intersects the output in the open band
$l+e<\theta_i<l+2e$; its norm enlargement does not change the derivative
weight $b$. Near contact, increasing the second height to $l+3e/2$
with other coordinates fixed stays in the output and in $\rho>0$:
$\partial_{\theta_2}\rho=1$ at contact and remains positive nearby.
At the endpoint the new and old functions agree on their actual cross
intersection. Second-variable unique continuation along this path
identifies the near-contact positive-side germs. The factors can be
exchanged. All choices have fixed positive margins and persist under
small smooth coefficient perturbations.
\end{proof}

\subsection{A compatible kernel at fixed separation}\label{subsec:fixed}
For a plane point $x_0$ and separation scale $s$, choose $\tau=1/100$.
Use first domain $B(x_0-\tau se_n,3\tau s)$ and observed radius
$(1-\alpha)\tau s$. In the second variable use
$s/2<|y-x_0|<2s$ and a small observed ball about $x_0-se_n$.
These domains are separated. Radial weights $|z-z_*|^{-1}$ have
normal Hessian $2|z-z_*|^{-3}$ and tangential value $-|z-z_*|^{-3}$,
so they are strict. A finite annular ball chain gives a fixed second
exponent $\gamma_2>0$. In the first ball, as $\alpha$ and a positive
output excess $cs$ decrease, $\gamma_1\to1$.
Choose them so $\gamma_1+\gamma_2>1$.
The cross lemma reaches the product of a ball about $x_0$ with a
narrower annulus about radius $s$.

Retain plane balls $|x-x_0|<cs/1000$ and smaller annuli, leaving
comparison margins inside the larger output products.
On overlaps the scales are comparable. Lowering the near-plane variable
by a small multiple of $c$ times the smaller scale stays in both larger
first balls and in their annular margins. Choose $\eta_0$ smaller than
these margins and than $\alpha\tau s/4$ over the compact scale range.
For each output, first compare with its own original prescription:
first-variable unique continuation from its observed ball identifies it
with $K_2$ on the connected lower part $x_n<-\eta_0$ of its output.
Second-variable continuation similarly identifies it with $K_1$ on the
lower part of its annulus. Those compact sublevels are known for small
$\lambda$. In an overlap of the same orientation, lower its near-plane
variable to this region. In an opposite-orientation overlap, lower both
variables by $\max(x_n,y_n,0)+2\eta_0$. Their difference is unchanged,
and the retained widths keep the paths inside both larger outputs.
Both germs then equal the actual common kernel. Unique continuation
identifies them on each overlap component. We have one kernel on the
coarse region
\begin{equation}\label{eq:seed}
 \min(x_n,y_n)<c_0|x-y|
\end{equation}
for a fixed $c_0>0$ on every fixed compact range of locations and positive
separations.

To extend further, use tubes $x=y+r_a(y)\omega$, $|\omega|=1$, where
\[
 t=M+y_n+e_0|y'|^2,\qquad r_a=a t^{1+\sigma},\qquad
 y_n\ge-B-1,\quad t\le T_0.
\]
Choose $M>B+3$ and small $\sigma>0$. Take $a_{\min}>0$ small enough to
pass the desired separations and $a_{\max}$ large enough to put its level
in \eqref{eq:seed}. Take $c_0a_{\min}T_0^\sigma>1$ and choose $e_0>0$
last. The base is compact; its bottom and $t=T_0$ faces are prescribed.
Every tube has positive minimum radius.

In the Euclidean limit set $\rho=(|x-y|^2-r_a^2)/2$, $b=\nabla r_a$, and
\[
 U_1=\omega+iv,\ v\perp\omega,\ |v|=1,\qquad
 U_2=\frac{\omega+b}{|\omega+b|^2}+iw,\quad
 w\perp(\omega+b),\ |w|=|\omega+b|^{-1}.
\]
Multiplying \eqref{eq:productcriterion} by $r_a^2$ gives
\begin{equation}\label{eq:tubetest}
 |\pi_{\omega^\perp}(U_1-U_2)|^2
       -(r_a\Hess r_a)(U_2,\overline{U_2}).
\end{equation}
Away from fixed neighborhoods of $b=-\omega,-2\omega$, the first term
is at least $c|b|^2$ on bounded $b$ ranges. Indeed its real part is
$-b_\perp/|\omega+b|^2$; if this vanishes, $b$ is parallel to $\omega$,
and the imaginary projected length equals $1/|\omega+b|$. The only
zeros are $b=0,-2\omega$. Near zero the real square controls
$|b_\perp|^2+O(|b|^3)$, and the imaginary length difference controls
$(\omega\cdot b)^2+O(|b|^3)$, uniformly over the permitted $v,w$.
Compactness gives the remaining lower bound.
Outside \eqref{eq:seed}, $r_a/t\le1/c_0$, so $b$ is bounded.
The two excluded neighborhoods are already prescribed: with
$\nabla t$ close to $e_n$,
$y_n/r_a+\omega_n\le(1+\sigma)|\nabla t|/|b|+\omega_n$
is small near $-\omega$ and negative near $-2\omega$.
Finally
\[
 r_a|\Hess r_a|\le C(\sigma+T_0e_0)|b|^2.
\]
Choose the exclusion sizes, then $\sigma$, then $T_0$, and lastly $e_0$
to absorb this error. Strictness persists for small spatial dilation.

Here is the compatibility argument used in this and the next continuation.
The prescribed region is upward closed in $a$. Let $a_*$ be a putative
last reached level. Each local crossing there uses compact data strictly
above $a_*+\eta_i$. A finite cover of the level gives outputs containing
base boxes times intervals $(a_*-h_i,a_*+h_i)$.
Every overlap has vertical segments to a common height above $a_*$,
where both germs equal the preceding function. Unique continuation on
neighborhoods of these segments identifies the germs on every component.
The finite cover crosses $a_*$, a contradiction. This argument uses bounds
on compact data above the last level, not a bound up to that level.
It supplies a single compatible kernel on every desired compact separated
region. To compare with the entire original cross, work in a larger region
and remove a small interior closed ball about a known $y$ from the
first-variable ball. The remaining domain is connected for $n\ge2$ and
contains a lower open ball where the kernels agree. First-variable unique
continuation identifies them; exchange variables for the other prescription.

\section{Shrinking spheres in a concave reference metric}\label{sec:spheres}

On the rescaled coordinate cylinder $Y=\{|y'|\le1,-1\le y_n\le1\}$ set
$t_0=y_n+|y'|^2$. Choose $L$ large and put
\[
 \gamma=e^{-2Lt_0}g_\lambda,\qquad
 t(y)=\int_{1/8}^{t_0(y)}e^{-2Ls}\dd s.
\]
The conformal connection formula gives
\begin{equation}\label{eq:depth}
 \Hess_\gamma t=e^{-2Lt_0}
   \{\Hess_{g_\lambda}t_0-L|\dd t_0|_{g_\lambda}^2g_\lambda\}
 \le-c\gamma
\end{equation}
after shortening the spatial range. In particular $\dd t\ne0$ and $t(0)<0$.
Choose $p$ large, then $R_*>0$ small, and define for $0<\delta\le1$
\begin{equation}\label{eq:radii}
 \ell_\delta(t)=\delta\log(1+e^{t/\delta}),\quad
 r_\delta=R_*\ell_\delta^p,\quad f_\delta=\log r_\delta,\quad
 \beta_\delta=|\dd f_\delta|_\gamma.
\end{equation}
The function $a(s)=e^s/\{(1+e^s)\log(1+e^s)\}$ is positive and decreasing;
it tends to one at $-\infty$ and is comparable to $(1+s)^{-1}$ at
$+\infty$, with $|a^{(j)}|/a^{j+1}$ bounded for each fixed $j$.
Thus \eqref{eq:depth} gives, uniformly in $\delta$,
\begin{equation}\label{eq:profile}
 \beta\ge\beta_0,\qquad\Hess_\gamma f\le-c\beta\gamma,\qquad
 |\nabla^jf|\le C_j\beta^j,\qquad r\beta^2\le\varepsilon.
\end{equation}
Here $\beta_0$ is made large by $p$, and the last bound follows from
$r\beta^2=R_*p^2\ell^{p-2}(\ell')^2|\dd t|^2$ for $p\ge2$.
The radii strictly increase with $\delta$, since
$\partial_\delta\ell_\delta=\log(1+e^s)-se^s/(1+e^s)>0$.
They tend uniformly to zero on compact subsets of $t<0$.

Multiplying the grounded equations by $e^{2Lt_0}$ gives
\begin{equation}\label{eq:conformal}
 \widetilde A_j=-\Delta_\gamma-2\ip{\dd\log\sigma_j}{\dd(\cdot)}_\gamma,
 \qquad\sigma_j=e^{(n-2)Lt_0/2}s_j.
\end{equation}
The symmetric density is $\sigma_j^2\vol_\gamma=e^{-2Lt_0}s_j^2\vol_g$.
Its reciprocal scalar change normalizes the delta exactly, so the kernels
$K_j$ and physical weighted fluxes are unchanged. We use these normalized
operators below and again denote them by $A_j$.

For a sphere $d_\gamma(x,y)=r(y)$ use
$\rho=(d_\gamma(x,y)^2-r(y)^2)/2$ and identify endpoints by parallel
transport. Put $b=\nabla r=r\nabla f$. The criterion vectors are the
$U_1,U_2$ above, and their denominators are bounded because
$|b|\le\varepsilon/\beta$. The short-geodesic Jacobi equation has
curvature coefficient $O(r^2)$ on parameter interval $[0,1]$.
Integrating twice with fixed endpoint values and using its index form
gives distance Hessian $|U_1-U_2|^2+O(r^2)$, with fixed geometric constants.
Consequently the strict expression is
\[
 |\pi_{\omega^\perp}(U_1-U_2)|^2-|b\cdot U_2|^2
             -r^2\Hess f(U_2,\overline{U_2})+O(r^2).
\]
The first term is at least $|b|^2-C|b|^3$. The orthogonal real and
imaginary vectors in $U_2$ have equal length $|\omega+b|^{-1}$, so
$|b\cdot U_2|^2\le |b|^2/|\omega+b|^2$. Concavity leaves
\begin{equation}\label{eq:spherestrict}
 cr^2\beta-Cr^2-Cr^3\beta^3>0
\end{equation}
for large $p$ and small $\varepsilon$. There is no principal-metric
mismatch term.

Fix any $\delta>0$. The spheres with $\delta\le\delta'\le1$ form a
compact off-diagonal cylinder. The initial level has positive minimum
radius. Choose $t_b>0$ with $3t_b<t(1/4)$. On $t>t_b$ all radii are
at least $R_*t_b^p$, so the fixed-separation kernel prescribes them.
On $y_n<-1/4$ the actual $K_1$ prescribes them. These regions cover
the base faces and are upward closed in $\delta'$. The vertical-box
compatibility argument in Section~\ref{subsec:fixed}, using
Lemma~\ref{lem:surface} and \eqref{eq:spherestrict}, extends one kernel
near the entire family. Distinct positive lower endpoints agree by the
same continuation from the common upper region. Thus there is one
compatible mixed kernel near every positive-radius sphere. This step
asserts no kernel bound uniform at radius zero.

Take a fixed cutoff $\chi$ compactly supported in $Y^\circ$, equal to
one near $0$, whose derivatives are supported in
\begin{equation}\label{eq:cutoffsupport}
 \{y_n<-1/4\}\cup\{t>t_b\}.
\end{equation}
For example multiply a cutoff rising between $y_n=-3/4,-1/2$ by one
falling between $t=2t_b,3t_b$. The depth bounds keep the lateral support
inside $|y'|<1$.

\section{The actual transfer density and its equation}\label{sec:trace}

For a first-side solution near the closed ball, define
\begin{equation}\label{eq:flux}
 \mathcal T_yv=\int_{\partial B_\gamma(y,r(y))}\!
 \sigma_1^2\{K(x,y)\partial_{\nu_\gamma}v(x)
                   -v(x)\partial_{\nu_\gamma}K(x,y)\}\dd S_\gamma(x).
\end{equation}
The first Dirichlet radial operator on this small ball moves to the smooth
kernel trace by sphere adjunction. Pulling back to the unit tangent sphere
$S_y$ with probability measure $\dd\omega_y$ gives a smooth density
\begin{equation}\label{eq:density}
 \mathcal T_yv=\int_{S_y}\psi(y,\omega)
            v(\exp_y^\gamma(r(y)\omega))\dd\omega_y.
\end{equation}
All surface Jacobians and weights are included in $\psi$.
For $y$ near a fixed center replace the moving boundary by a fixed slightly
larger surrounding surface lying in the kernel's open annulus. Green's
identity preserves the flux. The second operator can then be applied
under that fixed integral: $A_{2,y}\mathcal T_yv=0$.
On the lower known set the actual Green formula makes $\mathcal T_yv=v(y)$.
Unique continuation on connected $Y$ gives
\begin{equation}\label{eq:transfer}
 \mathcal T_y1=1,\qquad\mathcal T_yv_1=v_2(y)
\end{equation}
for every matching local pair. In particular $\psi-1$ has zero fiber mean.

We need the complete smooth parameter dependence of the small-ball radial
map, including its smoothing part. Here and below $\Psi^m$ means scalar
classical pseudodifferential operators of order $m$ on the unit sphere,
with local quantization
$(2\pi)^{-(n-1)}\int e^{i(\omega-z)\xi}a(\omega,\xi)u(z)\dd z\dd\xi$.

\begin{lemma}\label{lem:radial}
In smooth orthonormal center frames, the fixed-unit-ball Schr\"odinger
radial map is $D_{\rm Sch}(y,r)=B+r^2D_2(y,r)$, where $B$ is the Euclidean
radial DN map and $D_2$ is bounded and smooth in the complete $\Psi^1$
topology down to $r=0$. The grounded radial map is
\begin{equation}\label{eq:radial}
 D_{\rm gr}=B+r([B,m]-m)+r^2\widetilde D_2,
 \qquad m=\dd\log\sigma_1(y)\cdot\omega,
\end{equation}
with the same smooth remainder property.
\end{lemma}
\begin{proof}
The unit-ball metric is $g_{y,r}=r^{-2}\exp_y(r\cdot)^*\gamma=I+O(r^2)$.
Conjugation by $\sigma_1$ gives
$-\Delta_{g_{y,r}}+r^2Q(\exp_y(rz))$, with
$Q=\Delta_\gamma\sigma_1/\sigma_1$. The Dirichlet form is uniformly
coercive for small $r$ by the unit-ball Poincar\'e inequality.
Boundary integer regularity and the identity
$R_a-R_b=R_a(L_b-L_a)R_b$, followed by
$\partial R=-R(\partial L)R$, give smooth inverse dependence in every
integer norm on the fixed domain.

The boundary factor recursion \cite[Lemma 3.2 and Proposition 3.3]{Cekic}
has denominator twice the positive tangential root, uniformly bounded
below by $c|\xi|$. Its recursive symbols and all fixed parameter derivatives
therefore have uniform symbol bounds. To include the smoothing part,
write inward distance $s=h\tau$, $h=|\xi|^{-1}$. The conjugated boundary
operator times $h^2$ has expansion
$-\partial_\tau^2+\lambda^2+\sum_{j\ge1}h^jQ_j$, where
$\lambda\ge c$ and $Q_j$ have polynomial $\tau$ coefficients.
Start with $a_0=e^{-\lambda\tau}$; recursively solve
\[
 (-\partial_\tau^2+\lambda^2)a_j
       =-\sum_{l=1}^jQ_la_{j-l},\quad a_j(0)=0,\quad a_j\to0.
\]
Its forcing is a polynomial times $e^{-\lambda\tau}$.
Writing $a_j=e^{-\lambda\tau}P_j$, solve
$2\lambda P_j'-P_j''$ from highest polynomial degree downward and set
$P_j(0)=0$. Exponential decay bounds all fixed normal and parameter
seminorms. High-frequency cutoffs tending sufficiently rapidly to infinity
sum these symbols with arbitrarily small tails in each remainder topology.
Taylor remainders in $s=h\tau$ and cutoffs away from the boundary are
smoothing, uniformly with parameter derivatives. Finite chart patches
give a Poisson parametrix $H$ with exact trace after a smooth collar
correction and with $LH$ a smooth-kernel error.
The true Poisson operator is $H-RLH$. Integer estimates on $R$ give
bounds from every negative boundary Sobolev order to every interior
integer order for this correction and its parameter derivatives.
These imply all smooth-kernel seminorms. For each fixed integer order
$k\ge1$, extend in flattened charts by finite reflections with coefficients
$c_\ell$ satisfying $\sum_{\ell=1}^{k+1}c_\ell(-\ell)^m=1$ for
$0\le m<k$. The finite Vandermonde system matches the required normal
derivatives and gives a bounded $H^k$ extension. Apply Fourier
Cauchy--Schwarz to both kernel variables. Thus the true
radial map is a complete smooth $\Psi^1$ family, not just a pointwise
family modulo unspecified smoothing operators.

At $r=0$ the map is $B$. Metric, potential and radial vector have zero
first $r$ variation, so its first variation is zero. Taylor's integral
remainder proves the first assertion. On the boundary
$w=\sigma_1(\exp_y(r\omega))/\sigma_1(y)=1+rm+O(r^2)$ and
$\partial_\rho w=rm+O(r^2)$. Hence
$D_{\rm gr}=w^{-1}D_{\rm Sch}w-w^{-1}\partial_\rho w$ gives
\eqref{eq:radial}. Its linear term has angular order zero.
\end{proof}

Let $\nabla_i$ denote horizontal Levi--Civita differentiation on the unit
sphere bundle, preserving probability measure and acting covariantly on
all displayed frame indices. We use orthonormal center frames for the
fiber calculations; contracted formulas are independent of this choice. Let $A_i$ differentiate
the Euclidean harmonic extension in direction $i$, and take its fiber
adjoint $A_i^*$. On degree-$k$ harmonics $B=k$,
$[B,A_i^*]=A_i^*$ and $\sum_iA_i^*A_i^*=0$.
The center-variation Jacobi field with transported direction has value
$e_i$ and initial derivative zero; its endpoint is parallel $e_i+O(r^2)$.
Therefore the unit-ball velocity is $r^{-1}(e_i+O(r^2))+f_i\omega$.
Using the radial map for its normal component and taking the fiber adjoint
of the boundary differentiation gives
\begin{equation}\label{eq:center}
 C_i^*=r^{-1}A_i^*+f_iB+E_i+F_i,\qquad E_i\in r\Psi^1,\quad F_i\in\Psi^0.
\end{equation}
The order-zero terms include $\omega_i([B,m]-m)$ and
$rf_i([B,m]-m)$. The latter is bounded because $r\beta$ is bounded.
All fixed base derivatives of the normalized coefficients of order $j$
cost at most $C_j\beta^j$ by \eqref{eq:profile} and the complete parameter
bounds of Lemma~\ref{lem:radial}.

Write the second normalized equation locally as
$-\gamma^{ij}\nabla_i\nabla_j-b_2^i\nabla_i$; $b_2$ includes the fixed
coordinate/connection terms. Differentiating the transfer pairing gives
\begin{equation}\label{eq:Pequation}
 P\psi=0,\qquad
 P=r\{\gamma^{ij}(\nabla_i+C_i^*)(\nabla_j+C_j^*)
                         +b_2^i(\nabla_i+C_i^*)\}.
\end{equation}
This tests every smooth boundary trace: solve the first Dirichlet problem
on radius $r+\eta$ with a transported smooth trace $h$.
Its restriction to the inner sphere tends to $h$ in every fixed smooth
norm as $\eta\downarrow0$, by fixed-ball boundary regularity and smooth
inverse dependence. These solutions are homogeneous near the closed
inner ball. Thus the equation holds against every $h$, without imposing
an ill-posed smooth Cauchy problem.

Put $D_i=\nabla_i+f_iB$ and
\begin{equation}\label{eq:T}
 T=\sum_iA_i^*D_i,\qquad P=2T+U+J.
\end{equation}
The identities $D_i(r^{-1}A_j^*)=(r^{-1}A_j^*)D_i$ and
$\sum_i(r^{-1}A_i^*)^2=0$ cancel the apparent $r^{-2}$ terms.
Here $J$ is uniformly $\Psi^1$ with no base derivative, and $U$ has
base order $a=0,1,2$ with coefficients
\begin{equation}\label{eq:Uclass}
 U_a\in\varepsilon\beta^{-a}\Psi^{2-a}.
\end{equation}
A base derivative of order $j$ has the extra allowance $\beta^j$.
To see every term, put $R_i=r^{-1}A_i^*$ and $H_i=E_i+F_i$.
The remaining products are $rDD$, $rRE$, $rER$, $rDE$, $rED$,
$rRF$, $rFR$, $rDF$, $rFD$, $rHH$ and the smooth drift terms.
The $rRF,rFR$ and $rb_2R$ terms have bounded angular order one and enter
$J$. The others have respectively second-base coefficient $r\gamma^{ij}$,
first-base sizes $r\beta\Psi^1,r\Psi^0,r^2\Psi^1$, or zeroth-base sizes
$r\beta^2\Psi^2,r\Psi^2,r^2\beta\Psi^2,r\beta\Psi^1$ and lower orders.
All fit \eqref{eq:Uclass} using $r\beta^2\le\varepsilon$.
The second-base coefficient is real and independent of $\omega$.

\section{Angular coercivity and the joint density bound}\label{sec:coercivity}

For a fiber function use the harmonic decomposition
$\phi=\sum_{k\ge0}\phi_k$ and the spectral norm
$\norm{\phi}_s^2=\sum_k(1+k)^{2s}\norm{\phi_k}_{L^2(S_y)}^2$.
Harmonic polynomials are complete: Fischer decomposition gives polynomial
restrictions, and multivariate Bernstein approximation makes these dense
in continuous sphere functions, then in $L^2$. The integer norms through
order three agree with chart Sobolev norms by the sphere Laplacian
parametrix; its eigenvalue on degree $k$ is $k(k+n-2)$.
For a scalar angular $Q\in\Psi^m$, $m=1,2,3$, its two commutators with $B$
still have order $m$. Integer mapping and adjoint bounds give
\[
 \norm{\Pi_kQ\Pi_l}\le
 C\frac{\min((1+k)^m,(1+l)^m)}{(1+|k-l|)^2}.
\]
For $a,b\ge0$, $a+b=m$, the numerator is at most
$(1+k)^a(1+l)^b$. Cauchy--Schwarz with the summable discrete kernel
$(1+|k-l|)^{-2}$ therefore proves
\begin{equation}\label{eq:forms}
 |\ip{Qu}{v}|\le C\norm{u}_b\norm{v}_a.
\end{equation}
We use $(a,b)=(1/2,1/2),(1/2,3/2),(3/2,3/2)$; no fractional chart
interpolation is required.

For compact base support and zero fiber mean define
\begin{equation}\label{eq:Nnorm}
 \mathcal N(\phi)^2=\int_Y\{\norm{\nabla\phi}_0^2+
 \beta^2\norm{\phi}_1^2+\beta^{-1}\norm{\nabla\phi}_{1/2}^2+
 \beta\norm{\phi}_{3/2}^2\}\vol_\gamma.
\end{equation}
\begin{lemma}\label{lem:model}
Under \eqref{eq:profile}, for large $\beta_0$ there is $c>0$ independent
of the shrinking parameter such that
\[
 \norm{T\phi}^2-\norm{T^\flat\phi}^2\ge c\mathcal N(\phi)^2,
 \qquad T^\flat=\sum_iA_i(-\nabla_i+f_iB).
\]
\end{lemma}
\begin{proof}
At degree $k\ge1$, represent $v$ by a harmonic homogeneous polynomial
with Fischer norm $\norm{v}_F^2=\sum_\alpha\alpha!|v_\alpha|^2$.
Its relation to probability sphere norm is
$\norm{v}_F^2=N_k\norm{v}^2$,
$N_k=n(n+2)\cdots(n+2k-2)$. This follows by Gaussian integration by parts:
trace-free coefficient tensors eliminate pairings within either factor,
leaving precisely the Fischer sum. The radial moment satisfies
$\mathbb E|X|^{2k+2}=(n+2k)\mathbb E|X|^{2k}$.
Put $d=k+n/2-1$ and $a_i=\partial_{z_i}$.
For harmonic degree-$k$ polynomials $q_i$,
\begin{align}\label{eq:projection}
 \norm{\Pi_{k+1}\textstyle\sum_i z_iq_i}_F^2
 &=\sum_i\norm{q_i}_F^2+\sum_{ij}(a_jq_i,a_iq_j)_F
               -d^{-1}\norm{\textstyle\sum_i a_iq_i}_F^2,\\[-2pt]
 \Pi_{k+1}\textstyle\sum_i z_iq_i
 &=\sum_i z_iq_i-|z|^2\sum_i a_iq_i/(2d).\nonumber
\end{align}
Indeed $\Delta_z(|z|^2h)=4dh$ on degree $k-1$; multiplication and
differentiation are Fischer adjoints. The probability-sphere adjoint is
$A_i^*=(2d+2)\Pi_{k+1}z_i$ on degree $k$.

Write $D_i^\pm=\pm\nabla_i+kf_i$ and $C_{ij}=2k(\Hess_\gamma f)_{ij}$.
In the following calculation, $|\cdot|_F$ is the pointwise Fischer norm
and $\norm{\cdot}_F^2$ includes integration against $\vol_\gamma$.
Put
\[
 \begin{aligned}
 s^\pm(y)&=\bigl|\textstyle\sum_i a_iD_i^\pm v(y)\bigr|_F^2,
 &S^\pm&=\int_Ys^\pm\vol_\gamma,\\
 X^+&=\int_Y\sum_{ij}(a_jD_i^+v,a_iD_j^+v)_F\vol_\gamma,
 &C[av]&=\int_Y\sum_{ij}C_{ij}(a_iv,a_jv)_F\vol_\gamma.
 \end{aligned}
\]
Covariant integration by parts gives
\begin{align*}
 X^+&=S^--C[av]+O(k^2)\norm{v}_F^2,\\
 S^+&\le k\norm{D^-v}_F^2-C[av]+O(k^2)\norm{v}_F^2,\\
 \norm{D^-v}_F^2&=\norm{D^+v}_F^2+\int\tr C\,|v|_F^2.
\end{align*}
Here $D_j^-D_i^+-D_i^+D_j^-=-C_{ij}$ plus curvature.
A degree-$k$ rotation has norm $O(k)$, and contractions add one factor
$k$, giving the displayed error. Keep the input degree $k$ fixed when
commuting these base operators through polynomial derivatives.
Substituting into \eqref{eq:projection} gives
\begin{equation}\label{eq:raising}
 N_k(\norm{Tv}^2-\norm{T^\flat v}^2)
 =(2+2/d)\{d\norm{D^+v}_F^2-S^+-dC[av]\}
                  +2S^-+O(k^3)\norm{v}_F^2.
\end{equation}
Use the integrated bound on fraction $1-\theta$ of $S^+$ and on the
remaining fraction use the pointwise bound, with $\zeta=c_1/\beta$,
\[
 s^+\le(1-\zeta)(k+n-1)|D^+v|_F^2+
           \zeta\{2s^-+8k^3\beta^2|v|_F^2\}.
\]
Integrate this pointwise bound with its displayed weights before combining
with \eqref{eq:raising}. The first contraction bound is at most $k+n-1$;
the second follows from
$D^++D^-=2k\dd f$ and $\sum_i\norm{a_iv}_F^2=k\norm{v}_F^2$.
The coefficient retained on $|D^+v|^2$ is
\[
 (2+2/d)\{(n-2)/2-\theta(n-1)+\theta\zeta(k+n-1)\}.
\]
Choose $\theta>0$ small in terms of $n$. This is bounded below by
$c+c'k/\beta$. The $S^-$ coefficient is nonnegative for large $\beta_0$.
The remaining Hessian form is
\[
 -(2+2/d)\{(d-1+\theta)C[av]+(1-\theta)k\int\tr C|v|_F^2\},
\]
which controls $k^3\int\beta|v|_F^2$ for $n\ge3$, $k\ge1$.
Take $c_1$ small to absorb its loss and then $\beta_0$ large to absorb
curvature.

Both weighted cross terms have the needed sign. For $w\ge0$,
\[
 \int w|D^+v|^2=\int\{w|\nabla v|^2+wk^2\beta^2|v|^2
                               -k\Div(w\nabla f)|v|^2\}.
\]
For $w=1$, $\Delta f<0$. For $w=\beta^{-1}$,
$\Div(\nabla f/\beta)=\beta^{-1}\tr_{(\nabla f)^\perp}\Hess f\le-(n-1)c$.
Thus the coefficient $c+c'k/\beta$ controls all four terms of
\eqref{eq:Nnorm}. Divide by $N_k$ and sum. Horizontal derivatives preserve
degree and $T,T^\flat$ change it by one, so their distinct outputs are
orthogonal. First sum finitely many degrees. At each fixed $\delta>0$,
smoothness and dominated convergence pass to the full norms; constants
remain independent of $\delta$, without a uniform harmonic tail cutoff.
\end{proof}

\begin{lemma}\label{lem:Pcoercive}
For sufficiently small $\varepsilon$ and large $\beta_0$,
$\mathcal N(\phi)\le C\norm{P\phi}$ for every compact-base-supported
smooth $\phi$ with zero fiber mean.
\end{lemma}
\begin{proof}
Take full adjoints in $\vol_\gamma\dd\omega_y$. Direct degree adjunction
and $[B,A_i]=-A_i$ give $Z=T^\flat-T^\dagger=\sum_i f_iA_i$.
The exact squared-norm difference is
\begin{align}\label{eq:error}
 \norm{(2T+U)\phi}^2-\norm{(2T^\flat+U^\dagger)\phi}^2
 &=4(\norm{T\phi}^2-\norm{T^\flat\phi}^2)+\ip{\mathcal E\phi}{\phi},\\
 \mathcal E&=2[T^\dagger,U]+2[U^\dagger,T]+[U^\dagger,U]
                            -2\{UZ+(UZ)^\dagger\}.\nonumber
\end{align}
$T$ has (base order, angular order, size) $(1,1,1)$ or $(0,2,\beta)$;
$U$ has $(a,2-a,\varepsilon\beta^{-a})$, and $Z$ has $(0,1,\beta)$.
A scalar angular commutator lowers angular order by one; a base coefficient
derivative lowers base order by one and costs at most $\beta$.
The possible third-base commutators vanish because $U_2=r\gamma^{ij}$
is angular independent. For $[U^\dagger,U]$, write
$[U^\dagger-U,U]$; the real $U_2$ makes the first difference have base
order at most one, and its possible third-base coefficient also commutes
with $U_2$. Fourth-base terms are absent. The resulting types are
\[
 \varepsilon\beta^{-1}\Psi^1\partial_y^2,\qquad
 \varepsilon\Psi^2\partial_y,\qquad\varepsilon\beta\Psi^3.
\]
One base integration by parts in the first type and \eqref{eq:forms}
bound their forms by respectively
$C\varepsilon\int\beta^{-1}\norm{\nabla\phi}_{1/2}^2$,
$C\varepsilon\int\norm{\nabla\phi}_{1/2}\norm{\phi}_{3/2}$, and
$C\varepsilon\int\beta\norm{\phi}_{3/2}^2$.
For the middle type insert $\beta^{-1/2}$ and $\beta^{1/2}$ and use
Cauchy--Schwarz. Coordinate/horizontal derivative differences are fixed
angular order one and have the same bound. Thus
$|\ip{\mathcal E\phi}{\phi}|\le C\varepsilon\mathcal N(\phi)^2$.
Absorb this into Lemma~\ref{lem:model} and discard the nonnegative
comparison norm. Finally
$\norm{J\phi}\le C\norm{\phi}_1\le C\beta_0^{-1}\mathcal N(\phi)$
is absorbed after increasing $\beta_0$.
\end{proof}

Apply the last lemma to $\chi(\psi_\delta-1)$.
On the lower derivative support in \eqref{eq:cutoffsupport}, the density is
the ordinary small-ball Poisson evaluation density. Lemma~\ref{lem:radial}
makes it smooth down to $r=0$, with uniformly bounded angular norms and
first base derivative $O(1+r\beta)$. On $t>t_b$, radii have a positive
lower bound and all profile derivatives are uniformly bounded; the
fixed-separation kernel gives the same needed bounds. In particular,
angular norms and $r\nabla\psi_\delta$ are bounded on these supports.
Since $T1=0$, \eqref{eq:Uclass} gives a uniform bound for $P1$.
The second-base term of $[P,\chi]$ is
$r\gamma^{ij}(2\partial_i\chi\partial_j+\partial_i\partial_j\chi)$;
the other terms have at most one base derivative, with coefficients
$r\beta\Psi^1$, $r\Psi^0$, or bounded angular order one.
The support bounds therefore give
\[
 \norm{P[\chi(\psi_\delta-1)]}
 =\norm{-\chi P1+[P,\chi](\psi_\delta-1)}\le C.
\]
Coercivity proves the joint estimate
\begin{equation}\label{eq:joint}
 \sup_{0<\delta\le1}\norm{\chi\psi_\delta}_{L^2(Y\times S_y,
                    \vol_\gamma\dd\omega_y)}<\infty.
\end{equation}
This is not a pointwise-in-$y$ density bound. On a fixed small $Z$ with
$\chi=1$, $t<0$, \eqref{eq:transfer} gives for each matching pair
\[
 \norm{v_2-v_1}_{L^2(Z)}
 \le C_{v_1}\sup_Zr_\delta
          \norm{\psi_\delta}_{L^2(Z\times S_y)}\longrightarrow0.
\]
The left side is independent of $\delta$, so it is zero, and smoothness
gives equality on $Z$. The same $Z$ works for every pair; only the finite
constant $C_{v_1}$ changes. This completes Proposition~\ref{prop:local}.

\section{Global equality from all cap sources}\label{sec:global}

For real $F\in C_c^\infty(U_0)$ put $u_{j,F}=R_jF$. Define
\[
 \mathcal Q=\{p\in N^\circ:\ \text{there is an open }V\ni p
         \text{ such that }u_{1,F}=u_{2,F}\text{ on }V\text{ for every }F\}.
\]
The neighborhood precedes the universal source quantifier, so $\mathcal Q$
is open. Let $W$ be its component containing the connected cap $E$,
which belongs to $\mathcal Q$ by Lemma~\ref{lem:cap}.
Then $q_1=q_2$ on $W$: outside $\overline{U_0}$ choose a source with
nonzero value by Lemma~\ref{lem:jets} and subtract its two equations;
on $\overline{U_0}\Subset E$ the coefficients already agree.
For fixed $x\in W$, all-source equality makes
$G_1(x,\cdot)-G_2(x,\cdot)$ vanish against sources on $U_0$.
On $W$ its delta singularities cancel because the local operators are now
identical. Local regularity makes the difference a smooth homogeneous
solution even at $y=x$. Unique continuation on connected $W$ gives
$G_1=G_2$ on $W\times W$.

If $W$ is proper in $N^\circ$, it has an interior frontier.
In a compact coordinate neighborhood, choose a point of $W$ close enough
to its complement that its Euclidean distance to that closed complement
is smaller than its distance to the chart boundary. The inscribed ball
is contained in $W$ and touches its frontier at $p$. Its smooth boundary
is a known one-sided hypersurface near $p$; no regularity of the whole
frontier is required. Since $\overline{U_0}\Subset E$, $p$ is outside it.

Choose one real $F_0$ with $u_{1,F_0}(p)=1$. Equality from $W$ and
continuity give $u_{2,F_0}(p)=1$. Set $s_j=u_{j,F_0}$ and choose one
connected neighborhood $\Omega$ where they are positive, avoiding
$\overline{U_0}$. For every $F$ define $v_{j,F}=u_{j,F}/s_j$ there.
Their equations are \eqref{eq:ground} and their pairs match on the known
ball. Choose a finite set of $n$ sources whose first-side ratio gradients
span at $p$, using the arbitrary first jets; they remain spanning on a
smaller $\Omega$.

Proposition~\ref{prop:local} gives a single connected neighborhood $Z$
of $p$, meeting $W$, on which every ratio pair agrees. Subtract their
normalized equations. The principal parts cancel, leaving
\[
 \ip{\dd\log s_1-\dd\log s_2}{\dd v_{1,F}}_g=0.
\]
The spanning family gives $\dd\log(s_1/s_2)=0$ on $Z$.
Matching on $Z\cap W$ makes its constant zero. Hence $s_1=s_2$,
$q_1=\Delta_gs_1/s_1=\Delta_gs_2/s_2=q_2$, and all original source pairs
agree on $Z$. Thus $W\cup Z$ is a connected subset of $\mathcal Q$,
contradicting that $p$ is on the frontier of its component.

The interior of a connected smooth manifold with boundary is connected:
boundary half-balls meet only one interior component, so distinct interior
components would give disjoint relatively open pieces partitioning $N$.
Therefore $W=N^\circ$. The potentials agree on the original interior and
on its boundary by smooth continuity. This proves Theorem~\ref{thm:main}.

\section*{AI use}
This manuscript was prepared with assistance from OpenAI Codex in developing
and examining candidate arguments, checking source statements, and drafting
the exposition. The core continuation method is credited separately to its
mathematical source. The manuscript is submitted for independent mathematical
review.

\begin{thebibliography}{99}
\bibitem{Aronszajn} N. Aronszajn,
\emph{A unique continuation theorem for solutions of elliptic partial
differential equations or inequalities of second order},
Journ. de Math. \textbf{36} (1957), 235--249.
\url{https://sites.math.washington.edu/~blwilson/Nodal/Aronszajn.pdf}.
\bibitem{Cekic} M. Ceki\'c,
\emph{Calder\'on problem for Yang--Mills connections},
arXiv:1704.01362v2, revision 12 June 2018.
\url{https://arxiv.org/abs/1704.01362v2}.
\bibitem{HunterPDE} J. K. Hunter,
\emph{Notes on Partial Differential Equations}, revision 18 June 2014.
\url{https://www.math.ucdavis.edu/~hunter/pdes/pde_notes.pdf}.
\bibitem{HunterMeasure} J. K. Hunter,
\emph{Brief Notes on Measure Theory}, 27 January 2007.
\url{https://www.math.ucdavis.edu/~hunter/m201b/measure.pdf}.
\bibitem{MaSahooSalo} S. Ma, S. K. Sahoo and M. Salo,
\emph{The anisotropic Calder\'on problem at large fixed frequency on
manifolds with invertible ray transform},
J. London Math. Soc. \textbf{110} (2024), e13006.
\url{https://doi.org/10.1112/jlms.13006}.
\bibitem{OpenAIMetric} OpenAI,
\emph{Smooth anisotropic uniqueness in the Calder\'on problem from one
boundary patch}, preprint, 24 September 2026.
\href{https://github.com/openai/math/tree/adc7f1241b42e322a6451854ab7e4b4c146bf78a/preprints/Smooth-Anisotropic-Uniqueness-in-the-Calderon-Problem-from-One-Boundary-Patch-September-24-2026}{Pinned source and PDF (commit \texttt{adc7f1241b42}).}
\bibitem{OpenAIConnection} OpenAI,
\emph{Determination of a metric and a unitary connection from one boundary
patch}, preprint, 5 October 2026.
\href{https://github.com/openai/math/tree/fd4aeeb2ee4fc729c18d98444fed42fd0529eeeb/preprints/Determination-of-a-metric-and-a-unitary-connection-from-one-boundary-patch-October-5-2026}{Pinned source and PDF (commit \texttt{fd4aeeb2ee4f}).}
\bibitem{TaoFourier} T. Tao,
\emph{245C, Notes 2: The Fourier transform}, 6 April 2009,
Section 2, Theorem 21 and Corollary 22.
\url{https://terrytao.wordpress.com/2009/04/06/the-fourier-transform/}.
\bibitem{Taylor} M. Taylor,
\emph{Short Course on Pseudodifferential Operators}, author-hosted notes.
\url{https://mtaylor.web.unc.edu/wp-content/uploads/sites/16915/2022/05/psidolect.pdf}.
\end{thebibliography}
\end{document}
